\chapter{KẾ HOẠCH QUẢN LÝ DỰ ÁN}

\section{Tổng quan}

\subsection{Phạm vi và ước lượng dự án}

Nền tảng kết nối cộng đồng nhiếp ảnh phim với các dịch vụ phòng tối và phòng chụp là một hệ thống web hỗ trợ nhiếp ảnh gia tìm kiếm, đặt chỗ và thanh toán các dịch vụ liên quan đến phòng tối, studio và thiết bị. Hệ thốngổng đồng thời cung cấp các chức năng cho nhà cung cấp quản lý không gian, thiết bị, gói dịch vụ, workshop, đặt chỗ và hoạt động kinh doanh. Ngoài ra, hệ thống hỗ trợ quản trị viên giám sát người dùng, nhà cung cấp và các giao dịch trên nền tảng.

Dự án bao gồm các công việc chính sau:

\setlength{\LTleft}{0pt}
\setlength{\LTright}{0pt}
\renewcommand{\arraystretch}{1.3}

\begin{longtable}{|p{1.0cm}|p{6.0cm}|p{2.5cm}|p{2.5cm}|}
	\caption{\raggedright Cấu trúc phân rã công việc (WBS) và ước lượng công sức}
	\label{tab:wbs}\\
	\hline
	\textbf{WBS} & \textbf{Công việc} & \textbf{Độ phức tạp} &
	\textbf{ngày-người} \\
	\hline
	\endfirsthead
	
	\hline
	\textbf{WBS} & \textbf{Công việc} & \textbf{Độ phức tạp} &
	\textbf{ngày-người} \\
	\hline
	\endhead
	
	1 & Khởi tạo dự án & Trung bình & 4 \\
	\hline
	1.1 & Phân tích yêu cầu và hoàn thiện SRS & Trung bình & 2 \\
	\hline
	1.2 & Thiết lập môi trường, repository và công cụ phát triển & Đơn giản & 2 \\
	\hline
	2 & Phân tích và thiết kế & Trung bình & 8 \\
	\hline
	2.1 & Thiết kế cơ sở dữ liệu & Trung bình & 3 \\
	\hline
	2.2 & Thiết kế kiến trúc hệ thống và API & Trung bình & 2 \\
	\hline
	2.3 & Thiết kế giao diện người dùng & Trung bình & 3 \\
	\hline
	3 & Triển khai hệ thống & Phức tạp & 62 \\
	\hline
	3.1 & Xác thực và quản lý tài khoản & Đơn giản & 5 \\
	\hline
	3.2 & Quản lý hồ sơ người dùng và nhà cung cấp & Đơn giản & 4 \\
	\hline
	3.3 & Quản lý không gian sáng tạo & Trung bình & 8 \\
	\hline
	3.4 & Quản lý thiết bị & Trung bình & 7 \\
	\hline
	3.5 & Quản lý gói dịch vụ & Trung bình & 6 \\
	\hline
	3.6 & Quản lý đặt chỗ và xử lý trạng thái & Phức tạp & 12 \\
	\hline
	3.7 & Thanh toán và hóa đơn & Trung bình & 6 \\
	\hline
	3.8 & Workshop và đăng ký workshop & Trung bình & 5 \\
	\hline
	3.9 & Quản trị hệ thống và dashboard & Trung bình & 5 \\
	\hline
	3.10 & Thông báo và các chức năng hỗ trợ & Đơn giản & 4 \\
	\hline
	
	4 & Kiểm thử & Trung bình & 10 \\
	\hline
	4.1 & Kiểm thử đơn vị & Đơn giản & 3 \\
	\hline
	4.2 & Kiểm thử tích hợp & Trung bình & 4 \\
	\hline
	4.3 & Kiểm thử hệ thống & Trung bình & 3 \\
	\hline
	
	5 & Hoàn thiện và báo cáo & Trung bình & 8 \\
	\hline
	5.1 & Hoàn thiện và kiểm tra hệ thống & Đơn giản & 2 \\
	\hline
	5.2 & Hoàn thiện tài liệu và báo cáo & Trung bình & 6 \\
	\hline
	
\end{longtable}

\setlength{\LTleft}{0pt}
\setlength{\LTright}{0pt}
\renewcommand{\arraystretch}{1.3}

\subsection{Mục tiêu kiểm thử}

\begin{longtable}{|p{3.5cm}|p{6.5cm}|p{4cm}|}
	\caption{\raggedright Mục tiêu kiểm thử}
	\label{tab:testing-objectives}\\
	\hline
	\textbf{Giai đoạn kiểm thử} & \textbf{Phạm vi kiểm thử} & \textbf{Mục tiêu lỗi} \\
	\hline
	\endfirsthead
	
	\hline
	\textbf{Giai đoạn kiểm thử} & \textbf{Phạm vi kiểm thử} & \textbf{Mục tiêu lỗi} \\
	\hline
	\endhead
	
	Rà soát mã nguồn & Kiểm tra quy chuẩn lập trình và cấu trúc mã nguồn & < 10 \\
	\hline
	
	Kiểm thử đơn vị & Kiểm tra các chức năng và xử lý nghiệp vụ riêng lẻ & < 8 \\
	\hline
	
	Kiểm thử tích hợp & Kiểm tra sự kết nối giữa Frontend, Backend, API và cơ sở dữ liệu & < 5 \\
	\hline
	
	Kiểm thử hệ thống & Kiểm tra các luồng chức năng hoàn chỉnh của hệ thống & < 5 \\
	\hline
	
	Kiểm thử chấp nhận & Kiểm tra mức độ đáp ứng các yêu cầu chức năng của người dùng & < 3 \\
	\hline
	
\end{longtable}


\subsection{Các rủi ro dự án}
\setlength{\LTleft}{0pt}
\setlength{\LTright}{0pt}
\renewcommand{\arraystretch}{1.3}

\begin{longtable}{|p{0.6cm}|p{3.5cm}|p{3.5cm}|p{1.8cm}|p{4cm}|}

	\textbf{\#} &
	\textbf{Risk Description} &
	\textbf{Impact} &
	\textbf{Possibility} &
	\textbf{Response Plans} \\
	\hline
	\endhead
	
	1 &
	Timeline gấp (30 ngày) khiến các module core (đặt chỗ, thanh toán) chưa được hoàn thiện đúng hạn. &
	Sản phẩm demo không chạy được luồng chính, ảnh hưởng trực tiếp đến điểm đồ án. &
	High &
	Ưu tiên phát triển các tính năng lõi trước theo MoSCoW, cắt giảm ngay các tính năng phụ (AI, cộng đồng) nếu tiến độ trễ sau tuần 2. \\
	\hline
	
	2 &
	Xác nhận thanh toán qua VietQR được thực hiện thủ công, dễ xảy ra sai sót khi đối chiếu sao kê ngân hàng. &
	Đơn hàng bị xác nhận sai hoặc chậm trễ, gây trải nghiệm không tốt cho người dùng khi demo. &
	Medium &
	Thiết kế rõ trạng thái đơn hàng trung gian (\texttt{waiting\_confirmation}) và kiểm thử kỹ luồng đối chiếu trước khi tích hợp UI. \\
	\hline
	
	3 &
	Chức năng quản lý đặt chỗ và phân bổ tài nguyên (tránh trùng lịch phòng/thiết bị) hoạt động chưa chính xác. &
	Nhiếp ảnh gia có thể đặt trùng lịch với người khác, gây lỗi nghiêm trọng trong luồng chính của hệ thống. &
	High &
	Kiểm thử các trường hợp đặt trùng lịch, kiểm tra logic phân bổ tài nguyên ở Backend và xử lý lỗi rõ ràng trên Frontend. \\
	\hline
	
	4 &
	Tích hợp Frontend và Backend không đồng bộ về API hoặc cấu trúc dữ liệu. &
	Một số chức năng có thể không hiển thị hoặc hoạt động đúng trên giao diện. &
	Medium &
	Thống nhất API endpoint, request/response và kiểm thử API trước khi tích hợp giữa các module. \\
	\hline
	
	5 &
	Thành viên nhóm không hoàn thành task đúng tiến độ hoặc phát sinh xung đột khi merge code. &
	Ảnh hưởng đến tiến độ tích hợp và hoàn thiện hệ thống. &
	Medium &
	Phân chia task rõ ràng trên Jira, sử dụng Git branch riêng và thường xuyên cập nhật tiến độ của từng thành viên. \\
	\hline
	
	6 &
	Các lỗi phát sinh trong quá trình kiểm thử khiến một số chức năng chưa ổn định. &
	Ảnh hưởng đến chất lượng và khả năng hoàn thành bản demo cuối cùng. &
	Medium &
	Thực hiện kiểm thử sớm, ưu tiên sửa các lỗi ảnh hưởng đến chức năng chính và kiểm thử lại sau khi sửa lỗi. \\
	\hline
	
\end{longtable}


\section{Phương pháp quản lý dự án}

\subsection{Quy trình phát triển dự án}

Sau khi xem xét các mô hình phát triển phần mềm, nhóm lựa chọn Quy trình phát triển phần mềm lặp và gia tăng (Iterative and Incremental Software Development Process).

Theo mô hình này, hệ thống được chia thành nhiều phần nhỏ để phát triển từng bước. Sau mỗi vòng lặp, nhóm sẽ hoàn thành một số chức năng, kiểm thử và tiếp tục phát triển các chức năng tiếp theo.

Mô hình này phù hợp với đề tài vì hệ thống có nhiều module khác nhau và các chức năng có liên quan đến nhau. Việc phát triển từng phần giúp nhóm dễ theo dõi tiến độ, kiểm thử và sửa lỗi trong quá trình thực hiện.

Các lý do lựa chọn mô hình này gồm:
\begin{itemize}
	\item Phát triển các chức năng quan trọng trước.
	\item Dễ điều chỉnh khi có thay đổi yêu cầu.
	\item Có thể kiểm thử và sửa lỗi sau mỗi vòng lặp.
	\item Dễ theo dõi tiến độ và quản lý rủi ro.
	\item Có thể cải thiện hệ thống qua từng phiên bản.
\end{itemize}


\subsection{Quản lý chất lượng}

\subsubsection{Phòng ngừa lỗi}

Khi phát hiện lỗi, người phụ trách phải được thông báo ngay. Mỗi lỗi cần được đánh giá về mức độ nghiêm trọng và thời gian cần thiết để xử lý. Các lỗi quan trọng phải được ưu tiên giải quyết và có thời hạn xử lý cụ thể.

Nhóm cũng cần chủ động xác định các vấn đề có thể xảy ra trong quá trình phát triển để có phương án xử lý sớm.

\subsubsection{Review}

Quá trình review phải được thực hiện khách quan và không thiên vị thành viên nào. Khi phát hiện lỗi, người phụ trách phải được thông báo để xử lý.

Các lỗi phát hiện trong quá trình review cần được ghi nhận trên công cụ Bug Tracking, bao gồm mức độ ưu tiên. Người phụ trách có trách nhiệm đưa ra giải pháp và sửa lỗi trong thời gian phù hợp.

\subsubsection{Kiểm thử đơn vị}

Các test case phải được chuẩn bị đầy đủ và phù hợp với chức năng của hệ thống. Các trường hợp có thể xảy ra cần được xem xét để hạn chế bỏ sót lỗi.

Khi phát hiện lỗi trong quá trình Unit Testing, lỗi phải được ghi nhận trên công cụ Bug Tracking cùng với mức độ ưu tiên. Người phụ trách sẽ xử lý và sửa lỗi.

\subsubsection{Kiểm thử tích hợp}

Các test case phải bao quát việc kết nối giữa các module của hệ thống. Nhóm cần kiểm tra để đảm bảo các module hoạt động đúng khi kết hợp với nhau.

Các lỗi phát hiện trong quá trình Integration Testing phải được ghi nhận trên công cụ Bug Tracking và có mức độ ưu tiên cụ thể. Người phụ trách phải xử lý lỗi để đảm bảo các module hoạt động ổn định và chính xác khi tích hợp.

\subsubsection{Kiểm thử hệ thống}

System Testing yêu cầu các test case được thiết kế cẩn thận và phù hợp với yêu cầu cũng như thiết kế của hệ thống.

Các lỗi phát hiện trong quá trình System Testing phải được ghi nhận trên công cụ Bug Tracking, kèm theo các thông tin cần thiết như mức độ ưu tiên. Người phụ trách phải đưa ra giải pháp và xử lý lỗi trong thời gian phù hợp.

System Testing phải đảm bảo kiểm tra đầy đủ các chức năng của hệ thống, bao gồm cả việc kiểm tra tương tác giữa hệ thống và các hệ thống bên ngoài.

\subsection{Kế hoạch đào tạo}

\begin{table}[H]
	\centering
	\caption{Kế hoạch đào tạo}
	\renewcommand{\arraystretch}{1.3}
	
	\begin{tabular}{|p{4cm}|p{3.5cm}|p{2cm}|p{2.5cm}|}
		\hline
		\textbf{Nội dung đào tạo} & \textbf{Đối tượng} & \textbf{Thời gian} & \textbf{Thời lượng} \\ \hline
		
		Python (Backend) & Tất cả thành viên & Tuần 1 & 7 ngày \\ \hline
		
		HTML, CSS, JavaScript & Thành viên nhóm & Tuần 1 & 14 ngày \\ \hline
		
		Git, GitHub & Tất cả thành viên & Tuần 1 & 3 ngày \\ \hline
		
		PostgreSQL & Tất cả thành viên & Tuần 2 & 4 ngày \\ \hline
		
	\end{tabular}
\end{table}

\section{Các sản phẩm bàn giao của dự án}

\begin{table}[H]
	\centering
	\caption{Các sản phẩm bàn giao của dự án}
	\renewcommand{\arraystretch}{1.3}
	
	\begin{tabular}{|p{1cm}|p{3cm}|p{10cm}|}
		\hline
		\textbf{STT} & \textbf{Sản phẩm bàn giao} & \textbf{Ghi chú} \\ \hline
		
		1 & Sprint 1 &
		 Phân tích yêu cầu hệ thống \newline
		 Xác định Actor và Module \newline
		 Xác định công nghệ sử dụng \newline
		 Hoàn thành Report 1
		\\ \hline
		
		2 & Sprint 2 &
		 Hoàn thiện SRS \newline
		 Thiết kế cơ sở dữ liệu \newline
		 Thiết kế UI/UX \newline
		 Thiết lập project cơ bản \newline
		 Hoàn thành Report 2
		\\ \hline
		
		3 & Sprint 3 &
		 Implement Authentication \& User Management \newline
		 Implement Creative Space Management \newline
		 Implement Resource \& Equipment Management \newline
		 Implement Service Package Management \newline
		 Implement Reservation \& Resource Allocation \newline
		 Implement Service Session Management
		\\ \hline
		
		5 & Sprint 4 &
		 Hoàn thiện các chức năng còn lại \newline
		 Tích hợp Frontend và Backend \newline
		 Validation \& Constraint \newline
		 Integration Testing \& System Testing \newline
		 Sửa lỗi và Deployment \newline
		 Hoàn thành Final Report
		\\ \hline
		
	\end{tabular}
\end{table}


\section{Phân công trách nhiệm}

\begin{table}[H]
	\centering
	\caption{Ma trận phân công trách nhiệm}
	\renewcommand{\arraystretch}{1.3}
	
	\begin{tabular}{|p{4cm}|p{2.2cm}|p{2.2cm}|p{2cm}|p{2.2cm}|}
		\hline
		\textbf{Công việc} & \textbf{Thanh Nhàn} & \textbf{Minh Thư} & \textbf{Sơn Lâm} & \textbf{Quang Cường} \\ \hline
		
		Requirement Analysis & D & R & S & S \\ \hline
		
		Project Planning & D & S & & \\ \hline
		
		Database Design & D & D & S & S \\ \hline
		
		Backend Development & S & D & & \\ \hline
		
		Frontend Development & D & D & & \\ \hline
		
		Testing & S & D & & \\ \hline
		
		Documentation & D & R & S & S \\ \hline
		
	\end{tabular}
\end{table}

Trong đó:

\begin{itemize}
	\item D = Do: Người thực hiện chính.
	\item R = Review: Người kiểm tra.
	\item S = Support: Người hỗ trợ.
	\item Blank = Omitted: Không tham gia.
\end{itemize}


\section{Trao đổi thông tin trong dự án}

\begin{table}[H]
	\centering
	\caption{Trao đổi thông tin trong dự án}
	\renewcommand{\arraystretch}{1.3}
	
	\begin{tabular}{|p{3.5cm}|p{3.5cm}|p{3cm}|p{3cm}|}
		\hline
		\textbf{Trao đổi} & \textbf{Thành phần tham gia} & \textbf{Tần suất} & \textbf{Công cụ} \\ \hline
		
		Daily Meeting & Development Team & Daily & Google Meet \\ \hline
		
		Sprint Planning & All Members & Every Weeks & Google Meet \\ \hline
		
		Sprint Review & Supervisor & Every Sprint & Offline/Meet \\ \hline
		
		Issue Discussion & Related Members & As Needed & Google Meet \\ \hline
		
		Document Sharing & All Members & Continuous & Google Drive \\ \hline
		
		Source Code Review & Developers & Continuous & GitHub \\ \hline
		
	\end{tabular}
\end{table}


\section{Quản lý cấu hình}

\subsection{Quản lý tài liệu}

Tài liệu dự án được quản lý bằng:

\begin{itemize}
	\item Microsoft Word / LaTeX.
	\item Google Sheets.
	\item Google Drive.
\end{itemize}

\subsection{Quản lý mã nguồn}

Mã nguồn của dự án được quản lý bằng Git và GitHub.

\subsection{Công cụ và cơ sở hạ tầng}

\begin{table}[H]
	\centering
	\caption{Công cụ và cơ sở hạ tầng}
	\renewcommand{\arraystretch}{1.3}
	
	\begin{tabular}{|p{4cm}|p{9cm}|}
		\hline
		\textbf{Danh mục} & \textbf{Công cụ / Cơ sở hạ tầng} \\ \hline
		
		Technology & Python (Backend), ReactJS / Next.js (Web) \\ \hline
		
		Database & PostgreSQL \\ \hline
		
		IDEs / Editors & Visual Studio Code \\ \hline
		
		Diagramming & Draw.io \\ \hline
		
		Documentation & LaTeX, Microsoft Office, Google Docs/Sheets/Slides \\ \hline
		
		Version Control & GitHub, Google Drive \\ \hline
		
		Project Management & Jira \\ \hline
		
	\end{tabular}
\end{table}

\subsection{Chính sách quản lý phiên bản}

\begin{itemize}
	\item Mỗi tính năng được phát triển trên một nhánh riêng.
	\item Bắt buộc tạo Pull Request trước khi merge vào main.
	\item Cần review code trước khi approve merge.
\end{itemize}